\section*{अध्याय \ref{ch:b1:lab-techniques-1} --- प्रयोगशाला तकनीकें I: सुरक्षा, मापन और पृथक्करण}

\begin{solution}{exo:b1:lab-techniques-1:1}
\emph{ख़तरा}: H225 (अत्यंत ज्वलनशील द्रव, श्रेणी 2) इसकी माँग करता है, और
लेबल पर अधिक गंभीर शब्द होता है। H225 एक भौतिक \omterm{def:b1:lab-techniques-1:hazard}{संकट} है; H319
(आँखों में गंभीर क्षोभ) और H336 (उनींदापन या चक्कर) स्वास्थ्य
संकट हैं। सावधानियाँ: पास में कोई लौ या चिंगारी नहीं, बंद बोतल; चश्मा;
हवादार स्थान में या धूम-कोष्ठ में काम।
% ledger: ghs:acetone, hstat:H225, hstat:H319, hstat:H336
\end{solution}

\begin{solution}{exo:b1:lab-techniques-1:2}
(a) एक \omterm{def:b1:lab-techniques-1:hazard}{संकट}, अम्ल का एक गुण। (b) एक \omterm{def:b1:lab-techniques-1:hazard}{जोखिम}, छोटा क्योंकि
अनावरण छोटा है। (c) एक \omterm{def:b1:lab-techniques-1:hazard}{संकट}। (d) एक \omterm{def:b1:lab-techniques-1:hazard}{जोखिम}: वही \omterm{def:b1:lab-techniques-1:hazard}{संकट}, अधिक अनावरण
(वाष्प, छींटे) जब गरम और खुला हो।
\end{solution}

\begin{solution}{exo:b1:lab-techniques-1:3}
$R_f = 2.1/6.0 = 0.35$ और $4.2/6.0 = 0.70$। दूसरे धब्बे का
$R_f$ संदर्भ के बराबर है: नमूने में वह हो सकता है (यह संगत है), पर
एक ही निक्षालक पर समान $R_f$ सर्वसमता सिद्ध नहीं करते; दूसरा निक्षालक, या
नमूने और संदर्भ का ऐसा सह-धब्बा जो एक ही धब्बा बना रहे, निष्कर्ष को
दृढ़ करता है। 0.35 वाला धब्बा निश्चित रूप से कोई अन्य स्पीशीज़ है।
\end{solution}

\begin{solution}{exo:b1:lab-techniques-1:4}
अर्ध-चौड़ाई \qty{0.05}{mL}, आयताकार: $u = 0.05/\sqrt3 = \qty{0.029}{mL}$।
दो स्वतंत्र पठनों के अंतर के लिए,
$u = \sqrt{2} \times 0.029 = \qty{0.041}{mL}$।
\end{solution}

\begin{solution}{exo:b1:lab-techniques-1:5}
माध्य \qty{0.25100}{g}; विचलन $+2$, $-2$, $+5$, 0, $-5$
(\qty{e-4}{g} में), वर्गों का योग \qty{58e-8}{g^2}, $s = \sqrt{58 \times
10^{-8}/4} = \qty{3.8e-4}{g}$; $u = s/\sqrt5 = \qty{1.7e-4}{g}$;
$U = \qty{3.4e-4}{g}$: $m = (0.25100 \pm 0.00034)$~g, $k = 2$।
\end{solution}

\begin{solution}{exo:b1:lab-techniques-1:6}
एक भाग: $f = 1/(1 + 3 \times 30/50) = 0.357$, \qty{64.3}{\percent}
निष्कर्षित। \qty{15}{mL} के दो: $f = 1.9^{-2} = 0.277$,
\qty{72.3}{\percent}। \qty{10}{mL} के तीन: $f = 1.6^{-3} = 0.244$,
\qty{75.6}{\percent}।
\end{solution}

\begin{solution}{exo:b1:lab-techniques-1:7}
$z = 0.0012/\sqrt{0.0006^2 + 0.0002^2} = 0.0012/0.00063 = 1.9 \le 2$:
संगत, बस थोड़े से। \qty{0.0004}{mol/L} के साथ: $z = 0.0012/0.00045 = 2.7$,
संगत नहीं: अधिक परिशुद्ध \omterm{def:b1:titration-methods:titration}{अनुमापन} एक वास्तविक अंतर प्रकट करता है
(बनाने में त्रुटि, या \omterm{def:b1:titration-methods:titration}{अनुमापन} में)।
\end{solution}

\begin{solution}{exo:b1:lab-techniques-1:8}
Q: ठंडे में अल्प विलेय, गरम में अत्यंत विलेय। P ठोस को ठंडे की तुलना में गरम में
शायद ही बेहतर घोलता है (ठंडा करने पर कुछ क्रिस्टलित नहीं होता); R ठंडे में बहुत अधिक
घुला रखता है (बड़ी हानियाँ)। Q के साथ, \qty{5.0}{g} को कम से कम
$5.0/12 \times 100 = \qty{42}{mL}$ उबलता विलायक चाहिए, जो
ठंडे में $0.3 \times 0.417 = \qty{0.13}{g}$ घुला रखता है: अधिक से अधिक
\qty{4.9}{g}, \qty{97.5}{\percent}, पुनःप्राप्त हो सकता है (फ़िल्टर पर किसी हानि से
पहले)।
\end{solution}

\begin{solution}{exo:b1:lab-techniques-1:9}
एथेनॉल (\qty{20}{\celsius} पर $n_D = 1.3611$); जल (1.333) और
साइक्लोहेक्सेन (1.4266) अपवर्जित हैं। अपवर्तनांक ताप के साथ बदलता है,
इसलिए ताप के बिना कोई मान अर्थहीन है। जल और एथेनॉल के बीच का
1.350 का मान किसी मिश्रण का सुझाव देता है, जैसे जल युक्त एथेनॉल।
% ledger: nD:water, nD:ethanol, nD:cyclohexane
\end{solution}

\begin{solution}{exo:b1:lab-techniques-1:10}
$t = KV/V_a$ और $n$ भागों के साथ, $\ln f_n = -n\ln(1 + t/n)$; जब
$n \to \infty$, तब $n \ln(1 + t/n) \to t$ (क्योंकि $\ln(1 + \varepsilon) \sim
\varepsilon$), और $f_n$, $\mathrm{e}^{-t}$ की ओर घटता है (\cref{prop:b1:lab-techniques-1:multiple-extractions})।
यहाँ $t = 3 \times 30/50 = 1.8$, $\mathrm{e}^{-1.8} = 0.165$: अधिक से अधिक
\qty{83.5}{\percent}। तीन भाग ही \qty{75.6}{\percent} देते हैं, जो सीमा का
\qty{91}{\percent} है; उसके \qty{99}{\percent} तक पहुँचने के लिए 32
भाग चाहिए, हर एक \qty{1}{mL} से कम: दो-तीन भागों के बाद लाभ
परिश्रम के योग्य नहीं।
\end{solution}

\begin{solution}{exo:b1:lab-techniques-1:11}
$t = x - x_0$ के साथ: $\int_{-a}^{a} (a - |t|)/a^2 \,\mathrm{d}t =
2 \int_0^a (a - t)/a^2 \,\mathrm{d}t = 2 (a^2/2)/a^2 = 1$। सममिति से माध्य
$x_0$ है; प्रसरण $2 \int_0^a t^2 (a - t)/a^2
\,\mathrm{d}t = \frac{2}{a^2}\left(\frac{a^4}{3} - \frac{a^4}{4}\right) =
\frac{a^2}{6}$ है, इसलिए $u = a/\sqrt6 = 0.41\,a$, जबकि आयताकार स्थिति के लिए
$a/\sqrt3 = 0.58\,a$: यह जानना कि मध्य अधिक संभाव्य है,
अनिश्चितता घटा देता है।
\end{solution}

\begin{solution}{exo:b1:lab-techniques-1:12}
$c = m/(MV) = 0.25100/(204.2 \times 0.10000) = \qty{0.012292}{mol/L}$।
\omterm{def:b1:lab-techniques-1:uncertainty}{आपेक्षिक अनिश्चितताएँ}: द्रव्यमान $1.7 \times 10^{-4}/0.2510 = 6.8 \times
10^{-4}$, आयतन $0.06/100.00 = 6.0 \times 10^{-4}$; संयुक्त
$\sqrt{(6.8^2 + 6.0^2)} \times 10^{-4} = 9.1 \times 10^{-4}$, अर्थात्
\qty{0.091}{\percent}; $u(c) = \qty{1.1e-5}{mol/L}$, $U = \qty{2.2e-5}{mol/L}$:
$c = (0.012292 \pm 0.000022)$~mol/L, $k = 2$। तौल थोड़ा-सा
प्रभावी है; दोनों स्रोत एक ही आकार के हैं।
\end{solution}

\begin{solution}{pb:b1:lab-techniques-1:1}
\textbf{1.} GHS02 ज्वलनशील, GHS05 संक्षारक, GHS06 तीव्र विषाक्तता, GHS07
हानिकारक या क्षोभक।
\textbf{2.} सैलिसिलिक अम्ल: \emph{ख़तरा}, H318 (आँखों की गंभीर
क्षति, श्रेणी 1) के कारण। ऐस्पिरिन: \emph{चेतावनी} (केवल H302)।
\textbf{3.} \omterm{def:b1:lab-techniques-1:hazard}{संकट} निश्चित है, अनावरण नहीं: थोड़ा आयतन,
धूम-कोष्ठ में मापा और प्रयुक्त, दस्तानों, चश्मे और प्रयोगशाला कोट के साथ,
बोतल तुरंत बंद, पास में कोई लौ नहीं।
\textbf{4.} H314 (त्वचा की गंभीर जलन और आँखों की क्षति): चश्मा और दस्ताने।
H332 (साँस में जाने पर हानिकारक) और H226 (ज्वलनशील द्रव और वाष्प): धूम-कोष्ठ
और लौ की अनुपस्थिति।
\textbf{5.} कई मिनट तक जल से सावधानीपूर्वक धोइए, संभव हो तो
कॉन्टैक्ट लेंस निकालते हुए, और धोना जारी रखिए (P305+P351+P338); फिर
चिकित्सकीय परामर्श लीजिए।
\textbf{6.} उदासीनीकरण के बाद जलीय अपशिष्ट में, सिंक में
नहीं।
\textbf{7.} \ce{C7H6O3 + C4H6O3 -> C9H8O4 + C2H4O2}: C $7 + 4 = 9 + 2$,
H $6 + 6 = 8 + 4$, O $3 + 3 = 4 + 2$।
\textbf{8.} $2.00/138.0 = \qty{1.449e-2}{mol}$; अधिक से अधिक
$\num{1.449e-2} \times 180.0 = \qty{2.61}{g}$ ऐस्पिरिन।
\textbf{9.} अपरिष्कृत ठोस में अनभिकृत सैलिसिलिक अम्ल, एथेनोइक
अम्ल और जल है: उसका द्रव्यमान ऐस्पिरिन का द्रव्यमान नहीं है, और उसकी शुद्धता
अज्ञात है।
\textbf{10.} ठंडा करने के बाद जो घुला रहता है वह खो जाता है: ठंडे में छोटी
और गरम में बड़ी \omterm{def:b1:precipitation:solubility}{विलेयता} का अर्थ है कि उत्पाद गरम होने पर घुल जाता है और ठंडा होने पर
लगभग पूरा वापस आ जाता है।
\textbf{11.} हर अतिरिक्त मिलीलीटर ठंडा होने पर उत्पाद का अपना अंश
घुला रखता है।
\textbf{12.} धीमी वृद्धि बड़े, अधिक नियमित \omterm{def:b1:crystals-metals:crystal}{क्रिस्टल} देती है जो कम
मातृद्रव और अशुद्धियाँ फँसाते हैं; बर्फ़ \omterm{def:b1:precipitation:solubility}{विलेयता} घटाती है, और इसलिए
हानि भी।
\textbf{13.} $\qty{0.045}{L} \times \qty{3.3}{g/L} = \qty{0.15}{g}$।
\textbf{14.} $1.95/2.61 = \qty{75}{\percent}$ (\qty{74.7}{\percent})।
\textbf{15.} $807/6 = \qty{134.5}{\celsius}$।
\textbf{16.} विचलन $-0.5$, $0.5$, $-1.5$, $0.5$, $-0.5$, $1.5$; वर्गों का
योग 5.5; $s = \sqrt{5.5/5} = \qty{1.05}{\celsius}$।
\textbf{17.} $u_A = 1.05/\sqrt6 = \qty{0.43}{\celsius}$।
\textbf{18.} अर्ध-चौड़ाई \qty{0.5}{\celsius}: $0.5/\sqrt3 = \qty{0.29}{\celsius}$।
\textbf{19.} $1/\sqrt3 = \qty{0.58}{\celsius}$;
$u = \sqrt{0.428^2 + 0.289^2 + 0.577^2} = \sqrt{0.600} = \qty{0.77}{\celsius}$।
\textbf{20.} $U = 2 \times 0.775 = \qty{1.5}{\celsius}$:
$T = (134.5 \pm 1.5)$~\unit{\celsius}, $k = 2$।
\textbf{21.} वह नीचे और \qty{8}{\celsius} के परास में पिघलता है: वह
अशुद्ध है (सैलिसिलिक अम्ल, एथेनोइक अम्ल, जल)।
\textbf{22.} इकाई तक दिया गया मान $\pm\qty{0.5}{\celsius}$ के भीतर है:
$u_{\mathrm{ref}} = 0.5/\sqrt3 = \qty{0.29}{\celsius}$।
\textbf{23.} $2.3/6.0 = 0.38$ और $3.3/6.0 = 0.55$। अपरिष्कृत उत्पाद में
सैलिसिलिक अम्ल (संदर्भ के समान $R_f$) और ऐस्पिरिन थे; \omterm{def:b1:lab-techniques-1:recrystallisation}{पुनःक्रिस्टलन} के बाद
केवल ऐस्पिरिन का धब्बा बचता है: सैलिसिलिक अम्ल
हट गया है, कम से कम उस सीमा से नीचे तक जिसे प्लेट पकड़ती है।
\textbf{24.} ऐस्पिरिन पिघलते समय अपघटित होती है; धीरे गरम करने पर उसे
अपघटित होने का समय मिल जाता है, और बना मिश्रण नीचे पिघलता है। सारणीबद्ध मान
तीव्र तापन को संदर्भित करता है, जैसा बेंच पर होता है।
\textbf{25.} $z = |134.5 - 135|/\sqrt{0.600 + 0.083} = 0.5/0.827 =$
\textbf{0.60}: 2 से काफ़ी नीचे, मापे गए और सारणीबद्ध गलनांक
संगत हैं।
\end{solution}
